\section{Ideal polyhedra, angle cells, and a partition function}
\label{sec:idealcells}

The aim is to replace an integral over surfaces by ordinary integrals
over finitely many regions of angle assignments. Two issues have to
be settled: which assignments come from polyhedra, and what density
they carry. An \emph{ideal} vertex lies on the sphere at infinity of
hyperbolic three-space. It has no finite distance from an interior
point, but a polyhedron with ideal vertices can have finite volume.
Its boundary metric has a cusp at each vertex.

\subsection{The geometric parametrization}

In the polyhedral cell discussion take $n\ge4$; the three-ended sphere is
the zero-dimensional pants base case.

The \emph{dual graph} has a vertex for each face of the polyhedron
and an edge crossing each original edge. A dual face encircles one
original vertex. The orientation supplies a cyclic order of edges
at every dual vertex; a graph with these cyclic orders is a
\emph{ribbon graph}. Thickening its edges to narrow strips recovers
an oriented surface with boundary. Here the complementary faces
are labeled by the punctures. A simple circuit is a closed edge
path with no repeated vertices except its starting/ending vertex.

Rivin's intrinsic realization theorem identifies a complete finite-area
hyperbolic metric on an $n$-punctured sphere with the boundary metric of a
convex ideal polyhedron, uniquely up to ambient isometry. Flat doubles of
ideal polygons occur as limiting configurations~\cite{RivinIntrinsic}.
For oriented moduli one retains the orientation of the boundary and uses
orientation-preserving equivalences.

Fix a combinatorial polyhedron and write $\theta_e$ for its exterior dihedral
angles. In the dual graph, the ideal conditions are
\begin{align}
 &0<\theta_e<\pi,\label{eq:rivinangles1}\\
 &\sum_{e\in\partial f}\theta_e=2\pi
       &&\text{for every dual face }f,\label{eq:rivinangles2}\\
 &\sum_{e\in C}\theta_e>2\pi
       &&\text{for every simple dual circuit not bounding a face.}
       \label{eq:rivinangles3}
\end{align}
Here the dual graph is the planar, $3$-connected graph of a convex
polyhedron, with its oriented embedding fixed. With these hypotheses, these
are Rivin's characterization~\cite{RivinAngles}. Thus the admissible angle
assignments form an open convex polytope in an affine subspace. Generic
polyhedra have triangular faces; their dual maps are trivalent and have
$3n-6$ edges. The $n$ perimeter conditions are independent: their matrix
is the unsigned vertex--edge incidence matrix of the primal triangulation.
A row relation assigns numbers $c_v$ with $c_v+c_w=0$ across every edge;
a triangular face forces its three numbers to vanish, and connectedness
forces all of them to vanish. The dimension is therefore $2n-6$.

An edge of the dual graph has metric length $\theta_e$. The resulting metric
map has face perimeters $2\pi$ and no short nonfacial circuit. This is the
genus-zero instance of a $2\pi$-irreducible metric map. It gives the precise
combinatorial object behind the proposed sum over polyhedral types.

For example, an ideal tetrahedron has six edges and four ideal
vertices. At the regular tetrahedron each interior dihedral angle
is $\pi/3$, hence each exterior angle $\theta_e$ is $2\pi/3$.
The three angles around a dual face add to $2\pi$; a nonfacial
four-edge circuit has sum $8\pi/3>2\pi$. The four independent
linear perimeter equations leave $6-4=2$ free parameters, agreeing
with the real dimension of the four-punctured-sphere moduli space.
These are lengths of the \emph{dual metric graph}, measured in
radians; they are not the infinite lengths of the original ideal edges.

\subsection{A normalized cell integral}

A sum over combinatorial types must avoid counting the same labeled,
oriented object several times, and an integral on a constraint surface
must specify its measure. The first issue produces the automorphism
factor below. The second can already be seen in the toy constraint
$x+y=L$ with $x,y>0$:
\[
 \int_{x,y>0} f(x,y)\,\delta(x+y-L)\,\dd x\,\dd y
   =\int_0^L f(x,L-x)\,\dd x.
\]
Here $\delta$ is the Dirac delta, characterized by
$\int f(u)\delta(u-a)\,\dd u=f(a)$.
The measure on this line is $\dd x$, whereas its Euclidean arc length
is $\sqrt2\,\dd x$. Disintegration means choosing the measure on
each fiber so that integration over fibers and then over the constraint
value recovers the original ambient integral. This distinction is the
reason for retaining the determinant in the graph formula.

For a dual ribbon graph $G$ with labeled faces, define its constrained
Lebesgue measure by disintegrating $\prod_e\dd\theta_e$ with respect to the
face-perimeter map. Symbolically, at perimeters $\boldsymbol\alpha$,
\begin{equation}
 \dd\nu_G(\boldsymbol\alpha)=
 \left(\prod_{e\in E(G)}\dd\theta_e\right)
 \left(\prod_{i=1}^{n}
 \delta\left(\sum_{e\in\partial f_i}\theta_e-\alpha_i\right)\right).
 \label{eq:deltameasure}
\end{equation}
An edge is counted with its boundary multiplicity. This notation specifies
a disintegrated measure, rather than Euclidean hypersurface measure with an
unstated coarea factor.
Concretely, let $P$ be the perimeter matrix and split the edge coordinates
as $(u,v)$ so that the square $n\times n$ block $P_v$ is invertible.
On the fiber $P_u u+P_vv=\boldsymbol\alpha$, the measure is
\[
 \dd\nu_G=\frac{|\dd u_1\cdots\dd u_{2n-6}|}{|\det P_v|},
 \qquad v=P_v^{-1}(\boldsymbol\alpha-P_u u).
\]
Thus deleting $n$ dependent edge variables without this determinant would
change the normalization.

The ideal cell description is known to give a volume calculation, beyond
the possibility suggested by the realization theorem. Charbonnier--David--Eynard
compare the angle measure with the WP form through lambda lengths
\cite[Section 3.2]{CDE}. In Budd's normalization of metric-map volume,
\begin{equation}
 \boxed{
 V_{0,n}(0,\ldots,0)=
 2^{n-2}\sum_{[G]}\frac{1}{|\Aut(G)|}
 \int_{\mathcal P_G}\dd\nu_G(2\pi,\ldots,2\pi).}
 \label{eq:polyhedralpartition}
\end{equation}
Here $\mathcal P_G$ imposes positivity and irreducibility, $G$ ranges over
oriented trivalent maps with labeled faces, and automorphisms preserve this
data. Nontrivalent strata have zero volume. Budd gives this interpretation
and the normalization explicitly~\cite[Section 1.3]{Budd}.
Equivalently, his rooted-map measure gives each root weight
$1/(2|E(G)|)$; summing the $2|E(G)|/|\Aut(G)|$ rooted versions gives
the displayed unrooted weight. Reversing the cyclic orders is not an
allowed equivalence of oriented maps.

Thus there is a genuine partition function over polyhedral combinatorial
types. Formula~\eqref{eq:polyhedralpartition} is a finite cellwise integral
at fixed $n$; it is not a sum over geodesics. The factors in it encode three
choices that should never be suppressed: labels, orientation, and the
normalization of the constrained measure. Section~\ref{sec:tetrahedron}
checks them directly for $n=4$.

\subsection{The string-theoretic connection}

Saadi--Zwiebach's closed-string construction uses equal-perimeter polyhedral
diagrams and quadratic differentials~\cite{SaadiZwiebach}. Zwiebach's
minimal-area formulation imposes a lower bound $2\pi$ on nontrivial closed
curves and uses uniqueness to control the covering of moduli space
\cite{Zwiebach}. These give a concrete historical counterpart to the
suggestion that such cells should be used for integration.

Several identifications nevertheless require care. A quadratic differential is locally an expression $q=\varphi(w)\,\dd w^2$
whose transformation rule makes $|q|$ a well-defined flat metric away
from its zeros and poles. A Jenkins--Strebel differential is a special
one whose regular horizontal trajectories form cylinders of closed
curves. Their boundaries make a metric ribbon graph. Conversely a
suitable metric ribbon graph can be thickened by these flat cylinders
into a surface.
A hyperbolic polyhedron has an intrinsic hyperbolic boundary metric and also
extrinsic dual data. These constructions may share the same combinatorial
and length description without giving the same metric or the same conformal
coordinate on moduli. Likewise, a string amplitude includes its specified
integrand; a decomposition used for string integration does not automatically
select the WP measure. Equation~\eqref{eq:polyhedralpartition} uses an
explicit measure comparison, not just the common vocabulary of polyhedra.

